\documentclass[11pt,a4paper]{article}

% ------------------------------------------------------------------
% Paket dasar
% ------------------------------------------------------------------
\usepackage[utf8]{inputenc}
\usepackage[T1]{fontenc}
\usepackage{lmodern}
\usepackage[a4paper,margin=2.5cm,headheight=15pt]{geometry}
\usepackage{amsmath}
\usepackage{booktabs}
\usepackage{tabularx}
\usepackage{longtable}
\usepackage{array}
\usepackage{enumitem}
\usepackage{listings}
\usepackage{xcolor}
\usepackage[most]{tcolorbox}
\usepackage{titlesec}
\usepackage{fancyhdr}
\usepackage[hidelinks]{hyperref}

% ------------------------------------------------------------------
% Nama bagian dalam bahasa Indonesia
% ------------------------------------------------------------------
\renewcommand{\contentsname}{Daftar Isi}
\renewcommand{\tablename}{Tabel}
\renewcommand{\figurename}{Gambar}
\renewcommand{\refname}{Referensi}

% ------------------------------------------------------------------
% Tampilan
% ------------------------------------------------------------------
\definecolor{reportblue}{HTML}{1F4E79}
\definecolor{reportgray}{HTML}{F2F5F8}
\definecolor{codegray}{HTML}{F7F7F7}
\definecolor{codecomment}{HTML}{5A7A5A}

\titleformat{\section}
  {\Large\bfseries\color{reportblue}}
  {\thesection}{0.7em}{}
\titleformat{\subsection}
  {\large\bfseries\color{reportblue}}
  {\thesubsection}{0.6em}{}
\titleformat{\subsubsection}
  {\normalsize\bfseries}
  {\thesubsubsection}{0.5em}{}

\setlength{\parindent}{0pt}
\setlength{\parskip}{0.55em}
\setlength{\emergencystretch}{2em}
\renewcommand{\tabularxcolumn}[1]{>{\raggedright\arraybackslash}m{#1}}

\pagestyle{fancy}
\fancyhf{}
\fancyhead[L]{\small Laporan Mingguan}
\fancyhead[R]{\small Teknik Kompilator}
\fancyfoot[C]{\thepage}
\renewcommand{\headrulewidth}{0.4pt}
\renewcommand{\footrulewidth}{0pt}

\setlist[itemize]{topsep=0.35em,itemsep=0.2em,parsep=0.2em}
\setlist[enumerate]{topsep=0.35em,itemsep=0.25em,parsep=0.2em}

% ------------------------------------------------------------------
% Gaya listing kode
% ------------------------------------------------------------------
\lstdefinestyle{reportcode}{
  backgroundcolor=\color{codegray},
  basicstyle=\ttfamily\small,
  breaklines=true,
  breakatwhitespace=false,
  columns=fullflexible,
  keepspaces=true,
  frame=single,
  rulecolor=\color{black!20},
  framesep=6pt,
  xleftmargin=4pt,
  xrightmargin=4pt,
  showstringspaces=false,
  commentstyle=\color{codecomment}\itshape,
  keywordstyle=\color{reportblue}\bfseries,
  tabsize=2
}
\lstset{style=reportcode}

\newtcolorbox{infobox}[1][]{
  colback=reportgray,
  colframe=reportblue,
  boxrule=0.6pt,
  arc=2pt,
  left=8pt,
  right=8pt,
  top=6pt,
  bottom=6pt,
  #1
}

% ------------------------------------------------------------------
% Judul
% ------------------------------------------------------------------
\title{
  \vspace{-1.5cm}
  {\Large\scshape Teknik Kompilator}\\[0.35cm]
  {\Huge\bfseries Laporan Mingguan}\\[0.35cm]
  {\LARGE \textit{Function} dan \textit{Procedure}}
}
\author{}
\date{18 November 2026}

\begin{document}

% ==================================================================
% Halaman judul
% ==================================================================
\pagenumbering{roman}
\hypersetup{pageanchor=false}
\begin{titlepage}
\centering
\thispagestyle{empty}

\vspace*{2.2cm}
{\Large\scshape Teknik Kompilator\par}

\vspace{1.8cm}
{\Huge\bfseries Laporan Mingguan\par}
\vspace{0.7cm}
{\LARGE \textit{Function} dan \textit{Procedure}\par}

\vspace{2.2cm}
\begin{tabular}{>{\bfseries}rl}
Tanggal     & : 18 November 2026 \\[0.35cm]
\end{tabular}

\vfill
%{\large Laporan ini mencakup kemajuan, hambatan,\\
%dan rencana kerja untuk minggu berikutnya.\par}

\vspace{1.2cm}
\end{titlepage}
\setcounter{page}{2}
\hypersetup{pageanchor=true}

% ==================================================================
% Daftar isi
% ==================================================================
\tableofcontents
\newpage

\pagenumbering{arabic}

% ==================================================================
% 1. Kemajuan
% ==================================================================
\section{Kemajuan}

Target minggu ini adalah menambahkan dukungan \textit{function} dan
\textit{procedure} pada kompilator.

\begin{table}[ht]
\centering
\caption{Status target minggu ini}
\label{tab:target}
\begin{tabularx}{\linewidth}{@{}p{0.46\linewidth}p{0.16\linewidth}X@{}}
\toprule
\textbf{Target} & \textbf{Status} & \textbf{Berkas utama} \\
\midrule
\textit{Context checking} untuk function dan procedure
  & Selesai & \texttt{Context.java} \\
\textit{Code generation} aturan R untuk function dan procedure
  & Selesai & \texttt{Generate.java} \\
Penanganan hash table (\texttt{Bucket}) untuk setiap function/procedure
  & Selesai & \texttt{Bucket.java} \\
Penyesuaian aksi C dan grammar dengan dokumen course
  & Selesai & \texttt{Context.java}, \texttt{Parser.cup} \\
\bottomrule
\end{tabularx}
\end{table}

\subsection{\textit{Context checking} pada \texttt{Context.java}}

\begin{itemize}
  \item \texttt{C22} menyimpan \textbf{entry point} badan rutin pada
        \texttt{Bucket.orderNum} dan lexical level deklarasi pada
        \texttt{Bucket.lexicLev}.
  \item \texttt{C24}/\texttt{C26} menentukan kind rutin sebelum
        \texttt{R7}, sehingga metadata tidak lagi ditentukan secara
        tidak langsung dari \texttt{C0} atau \texttt{C5}.
  \item \texttt{C25} memberi offset parameter negatif terhadap frame
        base:
        \begin{itemize}
          \item function: parameter ke-\texttt{i} berada pada offset
                $-(i + 3)$;
          \item procedure: parameter ke-\texttt{i} berada pada offset
                $-(i + 2)$.
        \end{itemize}
  \item \texttt{C5} mencatat tipe parameter pada daftar parameter
        rutin, sehingga informasi parameter tetap tersedia setelah
        scope parameter ditutup oleh \texttt{C27}.
  \item \texttt{C30}--\texttt{C35} memakai \texttt{argCountStack}
        untuk menghitung parameter dan argumen. Ini juga mendukung
        pemanggilan bersarang.
  \item \texttt{C8} memakai nama pada \texttt{symbolStack}, bukan
        \texttt{currentStr}, karena pada pemanggilan function
        \texttt{C8} dieksekusi setelah argumen-argumen.
  \item \texttt{C16} tetap dipakai untuk assignment, sedangkan
        \texttt{C36} khusus untuk tipe hasil function.
\end{itemize}

\subsection{\textit{Code generation} pada \texttt{Generate.java}}

Aturan yang berkaitan dengan function dan procedure:

\begin{center}
\textbf{Tabel: Aturan code generation function dan procedure} \\[6pt]
\begin{tabular}{@{} l p{0.78\linewidth} @{}}
\hline
\textbf{Aturan} & \textbf{Fungsi} \\
\hline
\texttt{R6}  & Memindahkan hasil perhitungan, menghapus variabel, dan keluar dari expression scope. \\
\texttt{R42} & Return dari procedure. \\
\texttt{R43} & Return dari function. \\
\texttt{R44} & Memanggil procedure. \\
\texttt{R45} & Membentuk block mark untuk pemanggilan procedure. \\
\texttt{R46} & Membentuk block mark untuk pemanggilan function. \\
\texttt{R47} & Memanggil function. \\
\texttt{R48} & Mencatat/menyimpan argumen. Karena argumen sudah berada di stack, \texttt{R48} memperbarui counter argumen. \\
\texttt{R49} & Jika identifier adalah function maka persiapan pemanggilan function, else ambil alamat variabel. \\
\texttt{R50} & Jika identifier adalah function maka panggil function, else muat nilai variabel. \\
\hline
\end{tabular}
\end{center}



\subsection{Konvensi pemanggilan}

Urutan pemanggilan:

\begin{enumerate}
  \item Pemanggil mendorong slot \textit{return address} (placeholder).
  \item Argumen dievaluasi dan didorong.
  \item Argumen dibalik dengan rangkaian instruksi \texttt{FLIP} agar
        argumen pertama berada paling dekat dengan frame callee.
  \item Entry point badan rutin didorong, lalu instruksi \texttt{BR}.
  \item Callee menyiapkan frame melalui \texttt{R1}/\texttt{R2}.
  \item Saat return, \texttt{R5}/\texttt{R6} membongkar frame dan
        mengembalikan display.
  \item \texttt{R42}/\texttt{R43} mengambil return address relatif
        terhadap \textit{memory top}, lalu \texttt{BR}.
  \item Pemanggil membersihkan argumen:
        \begin{itemize}
          \item function: hasil dipindahkan ke slot return address,
                lalu argumen dibersihkan;
          \item procedure: argumen dan slot return address dibersihkan.
        \end{itemize}
\end{enumerate}

\begin{table}[ht]
\centering
\caption{Susunan frame saat pemanggilan function dan procedure}
\label{tab:frame}
\begin{tabularx}{\linewidth}{@{}lXX@{}}
\toprule
\textbf{Posisi} & \textbf{Function} & \textbf{Procedure} \\
\midrule
\texttt{R}                 & return address & return address \\
\texttt{R+1} \ldots \texttt{R+n} & argumen (setelah dibalik)
                           & argumen (setelah dibalik) \\
\texttt{R+n+1}             & slot hasil fungsi & --- \\
\texttt{R+n+2}             & display lama & display lama \\
\texttt{B}                 & \texttt{R+n+3} & \texttt{R+n+2} \\
\bottomrule
\end{tabularx}
\end{table}

Dengan susunan tersebut, parameter function ke-\texttt{i} berada pada
offset $-(i+3)$ dan parameter procedure ke-\texttt{i} berada pada
offset $-(i+2)$ relatif terhadap \texttt{B}.

\subsection{Hash table dan \texttt{Bucket.java}}

\texttt{Bucket} menyimpan metadata rutin:

\begin{itemize}
  \item \texttt{paramTypes}: daftar tipe parameter sesuai urutan
        deklarasi;
  \item \texttt{paramCount}: jumlah parameter;
  \item \texttt{lexicLev}: lexical level deklarasi;
  \item \texttt{orderNum}: alamat entry point badan rutin;
  \item \texttt{idKind}: \texttt{Bucket.FUNCTION} atau
        \texttt{Bucket.PROCEDURE}.
\end{itemize}

Dengan demikian, setiap function/procedure memiliki satu bucket yang
lengkap dan dapat dipakai untuk pemeriksaan argumen di titik
pemanggilan.

\subsection{Perubahan grammar pada \texttt{Parser.cup}}

\begin{table}[ht]
\centering
\caption{Aksi grammar untuk function, procedure, dan parameter}
\label{tab:grammar}
\begin{tabularx}{\linewidth}{@{}p{0.23\linewidth}X@{}}
\toprule
\textbf{Produksi} & \textbf{Aksi} \\
\midrule
Deklarasi function
  & \texttt{type FUNC IDENT C3 C26 R7 C22 C23 C5 C0 R2 funcBody} \\
Deklarasi procedure
  & \texttt{PROC IDENT C3 C24 R7 C22 C0 R1 procBody} \\
\texttt{funcBody} tanpa param.
  & \texttt{= expression C36 C11 C11 C2 R6 R43 R9 C7} \\
\texttt{funcBody} berparameter
  & \texttt{( C30 parameters C35 C1 ) = expression C36 C11 C11 C27 R6 R43 R9 C7} \\
\texttt{procBody} berparameter
  & \texttt{( C30 parameters C35 C1 ) scope C27 R5 R42 R9 C7} \\
\texttt{parameters}
  & \texttt{IDENT C3 C4 C25 : type C5 C11 C7 C34 moreParameters} \\
Panggil procedure (bare)
  & \texttt{C28 C29 R45 R44 C7} \\
Panggil procedure (argumen)
  & \texttt{( R45 C28 C30 arguments ) C32 R44 C7} \\
Panggil function (bare)
  & \texttt{C37 R49 C29 C8 R50 C7} \\
Panggil function (argumen)
  & \texttt{( C33 R46 C30 arguments ) C32 C8 R47 C7} \\
\texttt{arguments}
  & \texttt{expression C34 C31 C11 R48 moreArguments} \\
\bottomrule
\end{tabularx}
\end{table}

Bentuk pemanggilan dengan tanda kurung kosong (\texttt{f()} dan
\texttt{p()}) yang sebelumnya merupakan ekstensi implementasi sudah
dihapus agar grammar sesuai dengan dokumen course.

\section{Hambatan}

\begin{enumerate}
%  \item \textbf{\texttt{R7} dipakai untuk dua hal berbeda.}
%
%        Selain branch badan rutin, \texttt{R7} juga dipakai untuk
%        branch \texttt{else}. Pada implementasi sebelumnya, metadata
%        rutin disimpan saat \texttt{R7}. Solusinya, metadata dipindahkan
%        ke \texttt{C22} sehingga \texttt{R7} cukup membangkitkan
%        forward branch.

%  \item \textbf{Parameter tidak dialokasikan oleh grammar lama.}
%
%        Parameter sudah didorong pemanggil di bawah frame base.
%        Solusinya, \texttt{C4} mengisi LL/ON umum, lalu \texttt{C25}
%        memberi offset negatif sesuai posisi parameter dan jenis rutin.

%  \item \textbf{Scope parameter memakai \texttt{C27}, bukan \texttt{C2}.}
%
%        Badan function/procedure berparameter harus keluar dari scope
%        parameter sebelum \texttt{R6}/\texttt{R5}. Solusinya,
%        \texttt{C27} diimplementasikan sebagai penghapusan simbol pada
%        lexical level parameter dan penurunan lexical level.

  \item \textbf{Mesin tidak memiliki instruksi \texttt{CALL}/\texttt{RET}.}

        Pemanggilan dibangun dari \texttt{BR}, \texttt{PUSH},
        \texttt{LOAD}, \texttt{STORE}, \texttt{PUSHMT}, dan
        \texttt{POP}. Return address dan pembersihan argumen dihitung
        relatif terhadap memory top. Pembersihan argumen dan pemindahan
        hasil function dilakukan oleh pemanggil setelah callee return.

  \item \textbf{Informasi parameter hilang setelah scope keluar.}

        \texttt{C27} menghapus simbol parameter dari hash table.
        Solusinya, tipe dan jumlah parameter disalin ke bucket rutin
        saat deklarasi, sehingga tetap dapat dipakai untuk pemeriksaan
        argumen di titik pemanggilan.

  \item \textbf{Aksi \texttt{C37} sebelumnya belum ada.}

        Akibatnya identifier non-function pada \texttt{subsOrCall}
        tidak melewati \texttt{C20}. Array yang dipakai sebagai scalar
        lolos tanpa error. Sekarang \texttt{C37} memilih \texttt{C33}
        untuk function dan \texttt{C20} untuk selain function.

  \item \textbf{Keterbatasan bahasa/grammar.}

        Belum ada dukungan parameter array atau parameter
        \textit{by-reference}. Belum ada deklarasi maju
        (\textit{forward declaration}), sehingga mutual recursion
        belum bisa. Pemanggilan rutin nested setelah scope induknya
        selesai juga belum dijamin aman karena display parent sudah
        tidak aktif.
\end{enumerate}

\section{Rencana Minggu Depan}

\begin{enumerate}
  \item \textbf{Menambah pengujian regresi otomatis.}

        Membuat script yang menjalankan pengujian otomatis dengan
        menjalankan kode berdasarkan input yang diberikan dan
        membandingkan output. Kasus error juga
        dimasukkan, termasuk pemanggilan
        routine berparameter tanpa tanda kurung.

  \item \textbf{Memperbaiki penanganan error sintaksis.}

        Pada beberapa input dengan kesalahan sintaks, error recovery
        masih dapat menyebabkan \texttt{EmptyStackException} pada
        context checking. Rencananya, context checking dibuat lebih
        defensif terhadap \texttt{symbolStack} dan \texttt{typeStack}
        yang kosong.

%  \item \textbf{Dukungan parameter array dan \textit{by-reference}.}
%
%        Jika spesifikasi bahasa mengizinkan, menambahkan parameter
%        array dan mekanisme alamat. Perlu penyesuaian konvensi
%        pemanggilan dan offset parameter.

  \item \textbf{Perapian dan dokumentasi.}

        Menampilkan entry point dan daftar parameter pada opsi
        \texttt{-p}, serta membersihkan penggunaan API
        \textit{deprecated} seperti \texttt{new Integer} secara
        bertahap.

  \item \textbf{Pengujian kombinasi scope yang lebih dalam.}

        Menguji rutin nested beberapa level, rekursi pada rutin
        nested, dan kombinasi array lokal dengan parameter.
\end{enumerate}

\appendix
\section{Daftar Berkas yang Diubah}

\begin{table}[ht]
\centering
\caption{Berkas yang diubah atau ditambahkan}
\label{tab:files}
\begin{tabularx}{\linewidth}{@{}p{0.28\linewidth}X@{}}
\toprule
\textbf{Berkas} & \textbf{Keterangan} \\
\midrule
\texttt{Bucket.java}   & Metadata parameter, jumlah parameter, lexical level, dan entry point rutin. \\
\texttt{Context.java}  & Aksi \texttt{C22}--\texttt{C37} disesuaikan dengan \texttt{Semantics} kompilator kuliah. \\
\texttt{Generate.java} & Aturan \texttt{R6}, \texttt{R42}--\texttt{R50}, helper call/return. \\
\texttt{Parser.cup}    & Grammar function/procedure disesuaikan dengan \texttt{Grammar} kompilator kuliah. \\
\bottomrule
\end{tabularx}
\end{table}

\end{document}
